
\documentclass[11pt,a4paper]{report}

% ---- Idioma y codificación ----
% babel: el ÚLTIMO idioma de la lista es el principal (spanish). El inglés
% se carga porque el Abstract usa \begin{otherlanguage}{english}.
\usepackage[utf8]{inputenc}
\usepackage[T1]{fontenc}
\usepackage[english,spanish]{babel}

\usepackage[backend=biber,style=ieee]{biblatex}
\addbibresource{bib/refs.bib}

\usepackage{styles/mistilo}

% ---- Datos del documento ----
% Cada TFG nuevo solo cambia estas líneas.
%   \titulodoc[título corto para la cabecera]{título completo}
\titulodoc[Entrega de explotación de vulnerabilidad IFOR]{Entrega de explotación de vulnerabilidad IFOR}
\autordoc{Pablo Pastor Aguado}                                  % <- tu nombre
\autordoc{Santiago González Rodríguez}                                  % <- tu nombre
\autordoc{Hugo Adán de la Fuente}                                  % <- tu nombre
\autordoc{Álvaro López Primo}                                  % <- tu nombre
\titulacion{Ingeniería informática}                         % <- tu grado
\mencion{Tecnologías de la información y la comunicación}   % <- tu mención
% \fechadoc{Octubre de 2026}   % descomenta para fijar la fecha

% ---- Compilar solo un capítulo mientras trabajas ----
% En una memoria larga, compilar todo tarda. Descomenta para compilar
% únicamente el capítulo 4; los demás conservan su numeración (\include).
% Vuelve a comentarla al terminar o faltarán capítulos en el PDF.
% \includeonly{capitulos/04_implementacion}

\begin{document}

\makeportada                         % portada definida en mistilo.sty

% ---- Páginas previas: numeración romana (i, ii, iii...) ----
\pagenumbering{roman}
% \input inserta el fichero tal cual (sin salto de página extra).
% frontmatter/resumen.tex -- SOLO CONTENIDO
% \capitulosinnumero (definido en mistilo.sty) crea un capítulo sin
% número, lo añade al índice y ajusta la cabecera.
\capitulosinnumero{Resumen}
Este trabajo presenta un prototipo ficticio que reparte consultas sobre
datos enlazados comprimidos entre varios procesos. Se describe su
arquitectura, se implementa en Python y se evalúa con datos ilustrativos.
La memoria sirve como ejemplo de documento largo para practicar la
modularización de proyectos en \LaTeX.

% frontmatter/abstract.tex
% otherlanguage cambia el idioma SOLO dentro del entorno (separación
% silábica y textos automáticos). Requiere cargar english en babel.
\capitulosinnumero{Abstract}
\begin{otherlanguage}{english}
This work presents a fictitious prototype that distributes queries over
compressed linked data among several processes. It is only an example
document used to practise how to split a \LaTeX{} project into modules.
\end{otherlanguage}

% frontmatter/agradecimientos.tex
\capitulosinnumero{Agradecimientos}
A quienes han hecho posible este taller.


\renewcommand{\contentsname}{Índice}
\tableofcontents                     % índice general automático
\listoffigures                       % índice de figuras (usa \caption)
\listoftables                        % índice de tablas

% ---- Cuerpo: numeración normal (1, 2, 3...) ----
\cleardoublepage
\pagenumbering{arabic}

% \include empieza en página nueva y permite \includeonly.
% PARA AÑADIR UN CAPÍTULO NUEVO:
%   1) crea capitulos/07_xxx.tex que empiece por \chapter{...}
%   2) añade aquí:  \include{capitulos/07_xxx}

% ---- Apéndices ----
% \appendix hace que los capítulos siguientes se numeren A, B, C...
\appendix

\end{document}
